\section{R effect estimation: preserve the target before choosing the learner}
\label{exec20:rfx-scope}

A servicing team may ask which customers should receive an intervention, what
the intervention achieves on average, or how much additional credit should be
offered. These questions require different statistical objects. Let $X$ denote
information available before the decision, $A\in\{0,1\}$ a recorded intervention,
and $Y(a)$ its potential monetary outcome. The conditional effect
$\tau(x)=E\{Y(1)-Y(0)\mid X=x\}$, its population average, and its average among
customers actually treated are not interchangeable. For a continuous amount $W$,
the dose-response curve $\mu(w)=E\{Y(w)\}$ introduces another target. Identification
still requires the relevant consistency, exchangeability and support assumptions;
orthogonal nuisance estimation does not supply them. Forest and double-machine-
learning results impose their own sampling and regularity conditions.
\citep{dml2018,grf2018}

The R routes examined here have distinct strengths. The selected \texttt{grf}
interface exposes conditional forests and several average-effect targets;
\texttt{tmle} addresses observed binary point treatments and outcome missingness;
R \texttt{DoubleML} exposes explicit orthogonal-score models. \texttt{CausalGPS}
and \texttt{causaldrf} supply different paths to continuous-exposure analysis.
The proposed product should wrap their qualified statistical contracts before
considering a native port. Matching an algorithm name or programming language
does not establish matching targets or uncertainty.
\citep{exec13_grf_rpath,exec20_rfx_tmle,exec20_rfx_doubleml_r,
exec20_rfx_causalgps,exec20_rfx_causaldrf}

\subsection{An R score is a specification, not a Python alias}

The inspected R DoubleML revision is dated 1 September 2026 and declares
development version \texttt{1.0.1.9000}. Its R6 and \texttt{mlr3} implementation
must be assessed separately from the Python installation considered earlier.
The selected interactive-regression class requires one binary treatment. Write
$m(x)=P(A=1\mid X=x)$ and $g_a(x)=E(Y\mid A=a,X=x)$. Its ATE score uses
\begin{equation}
 \psi_b=g_1(X)-g_0(X)+\frac{A\{Y-g_1(X)\}}{m(X)}
 -\frac{(1-A)\{Y-g_0(X)\}}{1-m(X)},\qquad \psi_a=-1,
 \label{exec20:rfx-irm}
\end{equation}
and solves the empirical equation $\overline{\psi_a}\theta+
\overline{\psi_b}=0$. For ATTE, the selected code instead uses the control
outcome nuisance, a treated fraction calculated within each test fold, and
$\psi_a=-A/\widehat p$. It fits the treated-outcome nuisance for ATE and custom
scores, but does not require that fit for its built-in ATTE path. These choices
must survive the wrapper: changing the target cannot be implemented merely by
changing a returned field name. The propensity truncation option clips predictions;
it does not itself define a new trimmed population.
\citep{exec20_rfx_doubleml_r,dml2018}

The nuisance helper supplies explicit training and test indices to custom
\texttt{mlr3} resampling. Treatment-specific outcome fits restrict training
indices while retaining the requested prediction rows. This is distinct from
cross-validation used to tune a learner: the inspected tuning method defaults
to global tuning and has an explicit option for tuning separately within folds.
A proposed validation protocol must record both sets of folds and any use of
outcomes during tuning. Neither a default option nor the existence of custom
scores proves that an arbitrary nuisance learner or new score satisfies a theorem.
\citep{exec20_rfx_doubleml_r,dml2018}

For continuous $W$, the partialling-out score solves a different moment. Let
$e(X)=E(W\mid X)$, $\ell(X)=E(Y\mid X)$ and $v=W-e(X)$. Its population target,
when the denominator is positive, is
\begin{equation}
 \theta_{\mathrm{PLR}}=
 \frac{E[v\{Y-\ell(X)\}]}{E(v^2)}.
 \label{exec20:rfx-plr}
\end{equation}
Under the stipulated conditional model
$Y=a(X)+\beta(X)W+\epsilon$, with $E(\epsilon\mid W,X)=0$, conditioning shows
$E(vY\mid X)=\beta(X)\operatorname{Var}(W\mid X)$. Hence
\begin{equation}
 \theta_{\mathrm{PLR}}=
 \frac{E\{\operatorname{Var}(W\mid X)\beta(X)\}}
 {E\{\operatorname{Var}(W\mid X)\}},
 \label{exec20:rfx-weighted}
\end{equation}
rather than generally $E\{\beta(X)\}$. Two equally probable risk groups with
slopes 1 and 3 and treatment variances 1 and 4 give respectively $13/5$ and 2.
The difference is a population-target difference even with perfect nuisances.
These are exact authored laws, not measurements of package error.
\citep{exec20_rfx_doubleml_r,exec20_rfx_reference}

\subsection{Forests, clustering and the information needed for uncertainty}

In the selected \texttt{grf} R path, all-sample AIPW averaging uses scores and
sums weighted centered scores within declared clusters for its variance.
The continuous-treatment score learns conditional treatment variance with an
additional forest and divides a treatment residual by that variance. Positivity
of this denominator and stable estimation are substantive conditions. A returned
local slope cannot be relabelled as an arbitrary finite-dose effect. The overlap
target uses a separate residual-regression path; the documented method choice
does not bypass earlier guards. In particular, the inspected average-effect
function rejects TMLE with clusters or sample weights. Honesty in forest
construction, treatment residualization and independent-cluster inference solve
different problems.
\citep{exec13_grf_rpath,grf2018}

The selected R DoubleML base class draws partitions over cluster identifiers
and has separate one- and two-way cluster variance calculations. This does not
establish validity for arbitrary serial dependence or few calendar periods.
Moreover, its multiplier-bootstrap method explicitly stops for clustered data.
Its joint-interval method requires bootstrap output, whereas the ordinary
interval method uses a normal critical value with the reported standard error.
A wrapper must preserve this unsupported combination; silently substituting
row multipliers would change the sampling law. A cluster-aware pointwise
standard error and a jointly valid set of comparisons are different deliverables.
\citep{exec20_rfx_doubleml_r}

For servicing data, the request must therefore name the independent sampling
units, repeated-customer identifiers, common calendar shocks, and how entire
units are kept out of nuisance-training folds. The public prototype can compare
supported regimes with known generated truth before an expensive bank trial.
It must retain failed folds, absent treatment arms, unsupported uncertainty
requests and denominator failures. These are proposed admission requirements;
the source inspection does not execute either R engine.

\subsection{Missing outcomes do not make missing treatment labels observed}

The pinned \texttt{tmle} 2.1.1 help text describes outcome missingness, but its
\texttt{Delta} argument also mentions missing treatment assignment. The selected
function body provides a narrower operational contract: its argument verifier
requires no missing values in $A$ or $W$, and the public entrypoint stops when
verification fails. Earlier preprocessing can also fail on an incompletely
observed treatment vector. Passing \texttt{A=NULL}, or an all-zero vector,
selects a constant-treatment path by replacing $A$ with ones. The documentation
describes the no-variation case as mean-outcome estimation. This is not treatment
imputation, and the missingness regression still conditions on recorded $A$.
The proposal should reject arbitrary missing-$A$ admission for this inspected
path and retain the conflicting wording for clarification and native fixtures.
\citep{exec20_rfx_tmle}

There is a scientific reason to keep those fields separate. Consider an authored
law in which observation group $R$ and randomized treatment $A$ are independent
fair coins. When $R=1$, both worlds have $Y(a)=a$ and the treatment is recorded.
When $R=0$, the treatment label is erased: world I has $Y(a)=a$, while world II
has $Y(a)=1-a$. The observed tuple consists of $R$, the treatment label when
present, and $Y$. Both worlds give the same distribution of that tuple: each
observed group has a fair binary outcome. Nevertheless,
\begin{equation}
 E\{Y(1)-Y(0)\}_{\mathrm{I}}=1,
 \qquad E\{Y(1)-Y(0)\}_{\mathrm{II}}=0.
 \label{exec20:rfx-missing}
\end{equation}
This witness deliberately lacks treatment-observation positivity in group
$R=0$ and grants no restriction transporting effects from $R=1$. It does not
claim impossibility under every missingness model. At a stipulated USD 100
outcome scale, the two effects are USD 100 and USD 0 despite the identical
observations. More rows with the same missing information cannot distinguish
them; verified labels or justified additional restrictions can.
\citep{exec20_rfx_reference}

\subsection{Dose curves require units, support and complete transformations}

A conditional regression slope also differs from a dose curve in one homogeneous
group. Suppose incremental credit $W$ is measured in thousands of dollars,
is uniform on $[0,2]$, and the stipulated margin curve is
$Y(w)=20w^2-5w$ dollars. Direct polynomial integration gives
\begin{equation}
 \frac{\operatorname{Cov}(W,Y)}{\operatorname{Var}(W)}=35,
 \quad Y'(0)=-5,\quad Y(1)-Y(0)=15.
 \label{exec20:rfx-dose}
\end{equation}
Using the projection slope to value a USD 1,000 increase from zero would report
USD 35 in this law, whereas the stipulated contrast is USD 15. Expressing credit
in dollars changes the slope to $35/1000$; it does not change the finite monetary
contrast. This example illustrates target and unit errors, not bank profitability.
\citep{exec20_rfx_reference}

The inspected \texttt{CausalGPS} kernel path standardizes an exposure residual
and stores an estimated residual density. Its weighting helper divides an
estimated marginal exposure density by that stored quantity. Under a
location-scale model $W=e(X)+s(X)U$, with $s(X)>0$, the change of variables is
\begin{equation}
 f_{W\mid X}(w\mid x)=
 \frac{1}{s(x)}f_{U\mid X}\!\left(\frac{w-e(x)}{s(x)}\middle|x\right).
 \label{exec20:rfx-jacobian}
\end{equation}
The inspected kernel branch does not divide by $s(x)$ before returning its GPS.
With different conditional scales, replacing the original-unit density with
residual density changes relative weights. An exact two-row illustration with
scales 1 and 2, the same residual density and marginal numerator, gives weights
1 and 2 rather than 1 and 1. Applied to stipulated outcomes USD 0 and USD 100,
their weighted means are USD $200/3$ and USD 50. This is algebraic motivation
for a conformance fixture, not a measured error for every package output.
\citep{exec20_rfx_causalgps,exec20_rfx_reference}

Under a constant positive change of exposure units, both a correctly expressed
marginal density and conditional density acquire the same reciprocal unit
factor, so their ratio is invariant. A covariate-dependent omitted scale is a
different issue. A public test should vary units and conditional scales separately,
then trace the resulting weighting and response fits. The selected response
routine supplies parametric, semiparametric and nonparametric routes; it rejects
negative weights but changes an all-zero weight vector to equal weights.
The product should refuse an all-zero causal weighting request before that
fallback changes its meaning. Balance statistics and effective sample size are
useful diagnostics, not proof of exchangeability or calibrated uncertainty.
Trimming and complete-case filtering must also be reflected in the target
population and reported support.
\citep{exec20_rfx_causalgps}

The selected \texttt{causaldrf} \texttt{hi\_est} path follows a different design.
It fits a treatment model, fits an outcome regression on treatment and the
estimated score, then recalculates the score at each requested dose for every
row before averaging predicted outcomes. This counterfactual-score step must
be preserved. Its log and square-root branches evaluate densities on transformed
scales without applying the corresponding original-scale Jacobian in that
function. At a fixed positive dose those Jacobian factors are constant across
covariates, so their omission does not by itself prove a failure of a balancing
argument at that dose. It does change the density interpretation and the
representation supplied to a restricted response regression across doses.
Positive support, transformation rules and the whole fitted response surface
therefore need explicit comparison before adoption.
\citep{exec20_rfx_causaldrf}

The accompanying \texttt{get\_ci} helper receives a matrix of bootstrap
coefficients, forms linear or quadratic predictions, and sorts each dose column.
It neither generates the resamples nor refits the treatment model. Its separate
column order statistics are pointwise constructions, not a simultaneous curve
band. The selected lower index is $\lfloor0.025B\rfloor$, which is zero when
$B<40$; native edge cases need inspection rather than a claim that increasing
$B$ solves statistical validity. A proposed curve experiment must specify which
sampling units are resampled, which stages are refitted, and whether it evaluates
pointwise coverage, joint coverage, or a selected optimal-dose decision.
\citep{exec20_rfx_causaldrf,exec20_rfx_reference}

\subsection{A bounded reuse proposal and the remaining comparison}

For a first public prototype, the candidates are R DoubleML's supported
IRM/PLR score paths, \texttt{grf}'s target-specific forest paths, and
\texttt{tmle}'s observed point-treatment paths. Each remains conditional on
its complete statistical and software contract. The developed \texttt{grf}
review in Section~\ref{exec21:grfr-scope} follows nuisance provenance and
target-specific control flow. Section~\ref{sec:exec24-tmle-reuse} additionally
identifies a specific intermediate-variable risk-ratio bootstrap interval
mismatch, resample-count and missing-output conventions, and an unresolved
rights question. The affected interval is refused pending a reviewed correction;
these findings cannot be waived by the earlier point-treatment comparison.
Requests should carry target, units, treatment-label status, support, weighting,
folds and sampling units. Responses should preserve unsupported modes and
diagnostic failures. Native ports remain conditional on demonstrated semantic
and numerical parity; translating a wrapper does not reproduce a learner's
statistical contract.
\texttt{CausalGPS} remains a conditional candidate pending scale and downstream
weight checks, while \texttt{causaldrf} is a useful proposed parametric baseline
subject to transformation and full-pipeline uncertainty qualification.

The selected originals establish these specific contrasts, not a complete R
market census. Further candidates, including \texttt{bartCause},
\texttt{causalweight}, other continuous-exposure methods, and their aliases,
still require matched primary-source review. The Bayesian forest comparison
elsewhere in this chapter does not close these package-specific omissions. Package license declarations and
copyright notices do not settle all distribution or learner-dependency rights.
No R package has been executed in this source comparison. Public and generated
experiments can decide which wrappers merit a prototype; bank truth, deployment
acceptance and observed returns remain subsequent questions. The immediate
comparison must preserve the question being estimated before ranking a learner.
